Vector search returns nearby embeddings from a large database. GPU
parallelism can accelerate distance calculations, but raw high-dimensional
vectors can exceed device memory. Compression makes larger indexes
resident; its benefit then depends on how codes are read, candidates are
selected and survivors are refined. This thesis asks how JHQ's compact
hierarchy can become an efficient GPU search system while keeping those
operations correct and measuring the resulting quality--speed trade-off.

CAGRA navigates a graph; compressed IVF scans chosen lists \cite{cagra,johnson_gpu,ivfrabitq}. Stella shows why compact storage matters (\cref{sec:7-7}):
\[
 S_{\rm raw}=67.81\,\mathrm{GiB}
 >S_{\rm GPU}=\frac{32607}{1024}\,\mathrm{GiB}
 >S_{\rm resident}=20.44\,\mathrm{GiB}.
\]

JHQ screens with primary codes and refines survivors with residuals \cite{jhq}. We vary execution and survivor count on the same trained index, measuring recall, speed, memory and construction cost, including losses.
